% GENERATED FILE. Edit canonical lessons, then run npm run generate:books.
% canonical-source-hash: fnv1a64:200d10eb3ef5a653
% canonical-lessons: HI-C159-bhara, HI-C159-sona, HI-C159-jagna, HI-C159-daurna, HI-C159-rakhna

\chapter{Full, To sleep, To wake up, To run, To put}
\label{ch:hi-full-to-sleep-to-wake-up-to-run-to-put}
\hlchaptermodality{159}

\begin{chapteropening}
\textbf{By the end of this chapter:} \emph{I can say full, to sleep, to wake up, to run and to put.}

Complete the last lesson of chapter 159: I can say full, to sleep, to wake up, to run and to put.
\end{chapteropening}

\section[\texorpdfstring{\hi{भरा} (bharā)}{भरा (bharā)}]{\hi{भरा} (bharā) — full}

\label{lesson:HI-C159-bhara}



\begin{quote}
Before the new one: say the Hindi for low, then the Hindi for empty.
\end{quote}

\subsection*{You'll want to know: \hi{भरा}}
\textbf{\hi{भरा}} — \emph{bharā} — \textquotedblleft{}full\textquotedblright{}.

\begin{cousinweb}
\textbf{\hi{भरा}} (\emph{bharā}) is a tadbhava, a word inherited from Sanskrit and reshaped by sound change. It comes from \emph{bharati}, \textquotedblleft{}bears, fills\textquotedblright{} (CDIAL 9397). \emph{bharā}, \textquotedblleft{}full\textquotedblright{}, is the past form of \emph{bharnā}, \textquotedblleft{}to fill\textquotedblright{}.
\end{cousinweb}

\begin{grammarlens}[title={What you've built}]
One more word for this chapter.
\end{grammarlens}

\subsection*{Your turn}
\begin{itemize}
\raggedright
  \item \emph{Say it:} \emph{bharā}
  \item \emph{Say it:} \emph{bharā}, once more
  \item \emph{Say it:} say \emph{bharā}
  \item \emph{Recall:} say the Hindi for low, then the Hindi for empty, then say \emph{bharā} again
\end{itemize}

\begin{tcolorbox}[breakable,colback=teal!4,colframe=teal!35!black,title={Before you move on}]
What does \hi{भरा} mean? (\textquotedblleft{}Full\textquotedblright{}.) Say it once more.
\end{tcolorbox}
\section[\texorpdfstring{\hi{सोना} (sonā)}{सोना (sonā)}]{\hi{सोना} (sonā) — to sleep}

\label{lesson:HI-C159-sona}



\begin{quote}
Before the new one: say the Hindi for empty, then the Hindi for full.
\end{quote}

\subsection*{You'll want to know: \hi{सोना}}
\textbf{\hi{सोना}} — \emph{sonā} — \textquotedblleft{}to sleep\textquotedblright{}.

\begin{cousinweb}
\textbf{\hi{सोना}} (\emph{sonā}) is a tadbhava, a word inherited from Sanskrit and reshaped by sound change. It comes from \emph{svapati}, \textquotedblleft{}sleeps\textquotedblright{} (CDIAL 13902). It is a different word from \emph{sonā}, \textquotedblleft{}gold\textquotedblright{}, from Sanskrit \emph{suvarṇa} (CDIAL 13519).
\end{cousinweb}

\begin{grammarlens}[title={What you've built}]
One more word for this chapter.
\end{grammarlens}

\subsection*{Your turn}
\begin{itemize}
\raggedright
  \item \emph{Say it:} \emph{sonā}
  \item \emph{Say it:} \emph{sonā}, once more
  \item \emph{Say it:} say \emph{sonā}
  \item \emph{Recall:} say the Hindi for empty, then the Hindi for full, then say \emph{sonā} again
\end{itemize}

\begin{tcolorbox}[breakable,colback=teal!4,colframe=teal!35!black,title={Before you move on}]
What does \hi{सोना} mean? (\textquotedblleft{}To sleep\textquotedblright{}.) Say it once more.
\end{tcolorbox}
\section[\texorpdfstring{\hi{जागना} (jāgnā)}{जागना (jāgnā)}]{\hi{जागना} (jāgnā) — to wake up, to stay awake}

\label{lesson:HI-C159-jagna}



\begin{quote}
Before the new one: say the Hindi for full, then the Hindi for to sleep.
\end{quote}

\subsection*{You'll want to know: \hi{जागना}}
\textbf{\hi{जागना}} — \emph{jāgnā} — \textquotedblleft{}to wake up, to stay awake\textquotedblright{}.

\begin{cousinweb}
\textbf{\hi{जागना}} (\emph{jāgnā}) is a tadbhava, a word inherited from Sanskrit and reshaped by sound change. It comes from \emph{jāgrati}, \textquotedblleft{}are awake\textquotedblright{} (CDIAL 5175).
\end{cousinweb}

\begin{grammarlens}[title={What you've built}]
One more word for this chapter.
\end{grammarlens}

\subsection*{Your turn}
\begin{itemize}
\raggedright
  \item \emph{Say it:} \emph{jāgnā}
  \item \emph{Say it:} \emph{jāgnā}, once more
  \item \emph{Say it:} say \emph{jāgnā}
  \item \emph{Recall:} say the Hindi for full, then the Hindi for to sleep, then say \emph{jāgnā} again
\end{itemize}

\begin{tcolorbox}[breakable,colback=teal!4,colframe=teal!35!black,title={Before you move on}]
What does \hi{जागना} mean? (\textquotedblleft{}To wake up, to stay awake\textquotedblright{}.) Say it once more.
\end{tcolorbox}
\section[\texorpdfstring{\hi{दौड़ना} (dauṛnā)}{दौड़ना (dauṛnā)}]{\hi{दौड़ना} (dauṛnā) — to run}

\label{lesson:HI-C159-daurna}



\begin{quote}
Before the new one: say the Hindi for to sleep, then the Hindi for to wake up.
\end{quote}

\subsection*{You'll want to know: \hi{दौड़ना}}
\textbf{\hi{दौड़ना}} — \emph{dauṛnā} — \textquotedblleft{}to run\textquotedblright{}.

\begin{cousinweb}
\textbf{\hi{दौड़ना}} (\emph{dauṛnā}) is a tadbhava, a word inherited from Sanskrit and reshaped by sound change. It comes from \emph{dravati}, \textquotedblleft{}runs\textquotedblright{} (CDIAL 6624). CDIAL notes an added \emph{-ḍ-}.
\end{cousinweb}

\begin{grammarlens}[title={What you've built}]
One more word for this chapter.
\end{grammarlens}

\subsection*{Your turn}
\begin{itemize}
\raggedright
  \item \emph{Say it:} \emph{dauṛnā}
  \item \emph{Say it:} \emph{dauṛnā}, once more
  \item \emph{Say it:} say \emph{dauṛnā}
  \item \emph{Recall:} say the Hindi for to sleep, then the Hindi for to wake up, then say \emph{dauṛnā} again
\end{itemize}

\begin{tcolorbox}[breakable,colback=teal!4,colframe=teal!35!black,title={Before you move on}]
What does \hi{दौड़ना} mean? (\textquotedblleft{}To run\textquotedblright{}.) Say it once more.
\end{tcolorbox}
\section[\texorpdfstring{\hi{रखना} (rakhnā)}{रखना (rakhnā)}]{\hi{रखना} (rakhnā) — to put, to keep}

\label{lesson:HI-C159-rakhna}



\begin{quote}
Before the new one: say the Hindi for to wake up, then the Hindi for to run.
\end{quote}

\subsection*{You'll want to know: \hi{रखना}}
\textbf{\hi{रखना}} — \emph{rakhnā} — \textquotedblleft{}to put, to keep\textquotedblright{}.

\begin{cousinweb}
\textbf{\hi{रखना}} (\emph{rakhnā}) is a tadbhava, a word inherited from Sanskrit and reshaped by sound change. It comes from \emph{rakṣati}, \textquotedblleft{}guards\textquotedblright{} (CDIAL 10547).
\end{cousinweb}

\begin{grammarlens}[title={What you've built}]
One more word for this chapter.
\end{grammarlens}

\subsection*{Your turn}
\begin{itemize}
\raggedright
  \item \emph{Say it:} \emph{rakhnā}
  \item \emph{Say it:} \emph{rakhnā}, once more
  \item \emph{Say it:} say \emph{rakhnā}
  \item \emph{Recall:} say the Hindi for to wake up, then the Hindi for to run, then say \emph{rakhnā} again
\end{itemize}

\begin{tcolorbox}[breakable,colback=teal!4,colframe=teal!35!black,title={Before you move on}]
What does \hi{रखना} mean? (\textquotedblleft{}To put, to keep\textquotedblright{}.) Say it once more.
\end{tcolorbox}
